% ============================================================
% 國立中央大學 資訊管理學系 碩士論文 — 序言（preamble）
% 版面依郭鎧菘論文為主、余姵穎論文為輔（2026-10-19 由其 PDF 量測整理）：
%   A4；中文標楷體(DFKai-SB)＋英文 Times New Roman，內文 12pt；
%   章 20pt 置中「三、 系統設計與實作」、節 18pt「3-1」、小節 16pt「3-2-1」；
%   頁碼底部置中 10pt，前置頁羅馬數字、正文起阿拉伯數字；圖表全篇連續編號；參考文獻 APA。
% 編譯：XeLaTeX（見 build.ps1）。若系上有官方範本，以官方為準並調整本檔。
% ============================================================

% ---------- 字型 ----------
\usepackage{fontspec}
\setmainfont{Times New Roman}
\setsansfont{Arial}
\setmonofont{Consolas}
\setCJKmainfont{DFKai-SB}[AutoFakeBold=3, AutoFakeSlant=0.2]   % 標楷體
\setCJKsansfont{Microsoft JhengHei}                             % 微軟正黑體（圖表內文）
\setCJKmonofont{MingLiU}

% ---------- 版面 ----------
\usepackage[a4paper, left=3.0cm, right=2.5cm, top=2.8cm, bottom=3.0cm]{geometry}
\linespread{1.8}                 % 12pt 字約 26pt 行距，對應學長論文量測值（約 25.5–27pt）
\setlength{\parindent}{2em}
\setlength{\parskip}{0pt}

% ---------- 章節標題（20 / 18 / 16 / 14 pt） ----------
\ctexset{
  contentsname    = {目錄},
  listfigurename  = {圖目錄},
  listtablename   = {表目錄},
  chapter = {
    name       = {,、},
    number     = {\chinese{chapter}},
    format     = {\centering\fontsize{20}{30}\selectfont\bfseries},
    aftername  = {\hspace{0.5em}},
    beforeskip = {-1.7cm},
    afterskip  = {18pt},
    pagestyle  = {plain},
  },
  section = {
    name       = {,},
    format     = {\raggedright\fontsize{18}{27}\selectfont\bfseries},
    aftername  = {\hspace{0.8em}},
    beforeskip = {12pt},
    afterskip  = {6pt},
  },
  subsection = {
    name       = {,},
    format     = {\raggedright\fontsize{16}{24}\selectfont\bfseries},
    aftername  = {\hspace{0.8em}},
    beforeskip = {8pt},
    afterskip  = {4pt},
  },
  subsubsection = {
    name       = {,},
    format     = {\raggedright\fontsize{14}{21}\selectfont\bfseries},
    aftername  = {\hspace{0.8em}},
    beforeskip = {6pt},
    afterskip  = {3pt},
  },
}
% 節、小節編號：3-1、3-2-1（郭鎧菘論文格式）
\renewcommand{\thesection}{\arabic{chapter}-\arabic{section}}
\renewcommand{\thesubsection}{\thesection-\arabic{subsection}}
\renewcommand{\thesubsubsection}{\thesubsection-\arabic{subsubsection}}

% ---------- 頁碼：底部置中 10pt ----------
\usepackage{fancyhdr}
\pagestyle{fancy}
\fancyhf{}
\fancyfoot[C]{\fontsize{10}{12}\selectfont\thepage}
\renewcommand{\headrulewidth}{0pt}
\fancypagestyle{plain}{\fancyhf{}\fancyfoot[C]{\fontsize{10}{12}\selectfont\thepage}%
  \renewcommand{\headrulewidth}{0pt}}

% ---------- 圖表：全篇連續編號；圖標題在下、表標題在上 ----------
\usepackage{graphicx}
\graphicspath{{figures/}}
\usepackage{chngcntr}
\counterwithout{figure}{chapter}
\counterwithout{table}{chapter}
\usepackage{caption}
\captionsetup{labelsep=space, font={normalsize}, skip=6pt}
% 圖目錄／表目錄條目顯示「圖 1」「表 1」（郭鎧菘論文格式）：只在列表內覆寫條目編號，不影響一般目錄
\makeatletter
\renewcommand*\l@figure{\@dottedtocline{1}{0em}{3.4em}}
\renewcommand*\l@table{\@dottedtocline{1}{0em}{3.4em}}
\let\ori@listoffigures\listoffigures
\renewcommand\listoffigures{\begingroup
  \renewcommand*\numberline[1]{\hb@xt@\@tempdima{圖~##1\hfil}}%
  \ori@listoffigures\endgroup}
\let\ori@listoftables\listoftables
\renewcommand\listoftables{\begingroup
  \renewcommand*\numberline[1]{\hb@xt@\@tempdima{表~##1\hfil}}%
  \ori@listoftables\endgroup}
\makeatother
\captionsetup[table]{position=above}
\captionsetup[figure]{position=below}
\usepackage{booktabs, array, tabularx, longtable, multirow, makecell}
\usepackage{float}

% ---------- 數學 ----------
\usepackage{amsmath, amssymb, amsthm}

% ---------- 顏色與待確認標記 ----------
\usepackage{xcolor}
% \pending{...}：標示尚待確認、不可進入定稿的內容（紅字醒目；定稿前全文搜尋 pending 清零）
\newcommand{\pending}[1]{\textcolor{red}{【待確認：#1】}}
\newcommand{\todo}[1]{\textcolor{blue}{【待辦：#1】}}

% ---------- 參考文獻：APA（biblatex-apa + biber） ----------
\usepackage{csquotes}
\usepackage[backend=biber, style=apa, sorting=nyt]{biblatex}
\DeclareLanguageMapping{english}{english-apa}
\addbibresource{references.bib}

% ---------- 超連結（最後載入） ----------
\usepackage[hidelinks, bookmarksnumbered, unicode]{hyperref}

% ---------- 程式與路徑等寬字 ----------
\newcommand{\code}[1]{{\small\texttt{#1}}}
